\section{Arithmetic of the actual common theta tests}
\label{sec:theta-test-inputs}

\status{TT1--TT10 reconstruct the finite local calculations, the finite
unfolding support, the expansion-independent global sum, and its
coefficient-field descent. They prove the actual \texttt{SumOKK} and
\texttt{SumRatK} conclusions; the latter retains its literal
\texttt{FieldOK} premise. TT11--TT14 construct the elementary parts of
that field and reduce its jet part to a specific comparison of the
same moved vacuum. TT15--TT19 give the compact-period arithmetic,
normalization-ratio argument and compatible finite-field join, with
their remaining analytic and finite-Weil inputs displayed. They do
not infer \texttt{LinK}, \texttt{InjK}, \texttt{Formula2K}, a tower
formula or positive spectral eigenvalues from finite-sum arithmetic.}

Throughout, $\mathcal C$ denotes one literal \texttt{Choice R}, with
its fixed grid $G$, Weil datum $W$, row measures, normal datum $Z$,
frames, representatives and reference twist $t_0=\mu_0$. The
hypotheses called \texttt{HidaSetting} include the actual
\texttt{LevCh.Is}, the arithmetic host and prime exclusions, full
Weil types on both ends, standard additive character, standard
source lattice, and $\beta_X=1$. We retain these hypotheses whenever
we specialize a local calculation to $\mathcal C$. Put
$K=\operatorname{range}(\mathrm{embLL})\subset\mathbb C$.
Membership in $K$ is a field statement. No paragraph identifies an
integral Weyl lattice with its dual or replaces an integral span by
its rational span.

\paragraph{TT1: locally constant integrals are finite arithmetic sums.}
Let $k$ be a nonarchimedean local field, $V=k^d$, and let $\mu$ be
additive Haar measure. Suppose a compact open set $H$ contains zero
and $\mu(H)=1$. If $f:V\to K$ is locally constant with compact
support, then its complex integral belongs to $K$.

Here is the normalization argument, including the case where $H$
is a set rather than a subgroup. The nonzero fibers of $f$ form a
finite collection of compact open sets: compactness gives a finite
cover by neighborhoods on which $f$ is constant. For each such fiber,
a finite cover by translates of a sufficiently small additive open
subgroup supplies a subgroup preserving that fiber. Intersect these
subgroups, their negatives, and the corresponding subgroup for
$1_H$. The result is an open subgroup $U$ preserving $f$ and $1_H$.
Because $0\in H$, $U\subset H$, so $U$ is compact. The support of
$f$ and $H$ are finite unions of $U$-cosets. If $H$ has $N$ such
cosets, then $N>0$ and
\[
 \mu(U)=N^{-1},\qquad
 \int_V f\,d\mu=N^{-1}\sum_{a\in A}f(a),                 \tag{TT.1}
\]
where $A$ is any finite set of support representatives. Every
summand is in $K$, and $N^{-1}\in\mathbb Q\subset K$.
This proves integrability as well as field membership. In
particular it applies to the source's finite product
\texttt{Vars}, with each factor normalized by
$\mu_k(\mathcal O_k)=1$. No topological closure of $K$ is used.

\paragraph{TT2: singular matrix balls and rational multiplicative volumes.}
Let $\nu:k^\times\to\mathbb C^\times$ be ramified, meaning that
$\nu|_{\mathcal O_k^\times}$ is nontrivial. On a matrix ball
$B=X+cM_n(\mathcal O_k)$ containing a singular matrix $X$, the
integral
\[
 \int_{B\cap\mathrm{GL}_n(k)}
       \nu(\det g)|\det g|^s\,dg=0                         \tag{TT.2}
\]
holds with the source's left Haar measure and totalized integral.
If $n=0$ there is no singular $X$, so this case does not arise.
For $n>0$ choose a nonzero left-kernel row $w_0X=0$, scale it to
$w\in\mathcal O_k^n$ with some $w_i=1$, and choose a unit $u$
with $\nu(u)\ne1$. Then
\[
 M=I+(u-1)e_iw,\quad \det M=u,\quad MX=X,\quad
 M^{-1}=I-u^{-1}(u-1)e_iw.                                 \tag{TT.3}
\]
Both $M$ and $M^{-1}$ are integral. Thus left multiplication by
$M$ preserves $B\cap\mathrm{GL}_n(k)$ and its measure, while
multiplying the integrand by $\nu(u)$. Subtracting proves (TT.2).
For a genuinely convergent integral this is ordinary change of
variables. The same equality is valid for the totalized integral;
it is not an assertion of convergence.

Conversely, suppose every matrix in a compact open ball $B$ is
invertible. Let $J\subset\mathrm{GL}_n(\mathcal O_k)$ be open
with Haar volume one, and assume $\nu$ is trivial on
$1+\pi^e\mathcal O_k$. There is a compact open subgroup
\[
 U\subset J\cap(1+\pi^eM_n(\mathcal O_k))
\]
whose right action preserves $B$. To obtain it, cover the compact
set $B$ by finitely many matrix neighborhoods contained in $B$ and
use continuity of multiplication; perform the same construction
for the complement near the compact set in question. Equivalently,
a sufficiently deep principal subgroup preserves each ball by the
bound on $B^{-1}$. The determinant character and determinant norm
are constant on each right $U$-coset. The left cosets $gU$ have
volume $\mu(U)=[J:U]^{-1}\in\mathbb Q$. Hence, whenever
$\nu(\det g)|\det g|^s\in K$, the integral over $B$ is a finite
sum of elements of $K$ times this rational volume. Notice that the
normalization concerns $J$, not all of $\mathrm{GL}_n(k)$.

\paragraph{TT3: the actual frame-orbit integral has values in the field.}
Let $\Psi$ be supported on the matrix lattice $\operatorname{Mat}
(\Lambda)$ and invariant under translation by
$\pi^e\operatorname{Mat}(\Lambda)$, where $\Lambda$ is a diagonal
full lattice and $e\geq1$. Suppose its values, the values of the
ramified character $\nu$, and every half-integral determinant norm
occurring below lie in $K$. For the fixed full-rank matrix
$Y=Y_0h$ in \texttt{smallF}, consider
\[
 \mathcal F(Y)=\int_{\mathrm{GL}_a(k)}
       \Psi(gY)\nu(\det g)|\det g|^{D/2}\,dg.              \tag{TT.4}
\]
Its value is in $K$, with the source's normalization
$\operatorname{vol}(J)=1$ and literal conductor bound
$\operatorname{CondLE}(\pi,e,\nu)$.

Indeed $Y$ has a right inverse. The implication
$\Psi(gY)\ne0\Rightarrow gY\in\operatorname{Mat}(\Lambda)$
therefore bounds every entry of $g$. Multiplication by $Y$ also
shows that $g\mapsto\Psi(gY)$ is invariant under addition of
$\pi^NM_a(\mathcal O_k)$ for some $N$. Thus its support is a
finite union of such matrix balls, with a constant coefficient
$\Psi(gY)$ on each. A ball containing a singular matrix contributes
zero by TT2. On every remaining ball TT2 gives a value in $K$.
Summing proves (TT.4) when the original integrand is integrable.
If it is not integrable, its source integral is defined to be zero,
which is also in $K$. This last branch is essential: TT3 proves the
arithmetic of \texttt{smallF}, not its analytic absolute convergence.

\paragraph{TT4: finite Fourier flips preserve the specified arithmetic.}
Write $q=|\mathcal O_k/\pi|$. For the source's block lattice $L$
and complementary block lattice $L'$, finite Fourier exchange has
coefficients
\[
 q^{-e r}\,\psi(\langle X,Y\rangle),                       \tag{TT.5}
\]
where $r$ is the number of transformed entries and the representatives
range over the prescribed finite quotient. The inverse transform uses
the opposite character and its prescribed normalization. If $K$
contains the required additive-character values, both transforms
preserve $K$-valued functions. The identities governing their support
and period follow from
\[
 \sum_{x\in A}\chi(x)=
 \begin{cases}|A|&\chi=1,\\0&\chi\ne1,\end{cases}          \tag{TT.6}
\]
for a finite additive group $A$. In the second case translate by an
$x_0$ with $\chi(x_0)\ne1$ and cancel. Applying (TT.6) entry by
entry proves the exact exchange of a lattice and its annihilator;
there is no analytic Fourier-limit step.

The middle operation \texttt{phiA} is zero off its prescribed
primitive-frame lattice. On that lattice it is a determinant
character/norm factor times (TT.4), so it too is $K$-valued. It has
the claimed depth for the following explicit reason. If $X$ is a
primitive frame and $Z\in\pi^e\operatorname{Mat}(L')$, extend
$X$ to an integral basis of $L'$. In that basis the rows of $Z$ are
$\pi^e$-divisible; fill the remaining rows by zero. The resulting
matrix $I+A$, with $A\in\pi^eM(\mathcal O_k)$, is invertible
and gives $X+Z=X(I+A)$. Transporting back gives a lattice
congruence element $\kappa$. Its inverse transpose preserves the
dual lattice, its determinant is in $1+\pi^e\mathcal O_k$, and
the defining change of frame leaves \texttt{smallF} unchanged.
The conductor bound kills the determinant character and its norm
is one. Nonprimitive frames remain nonprimitive after this small
translation. These arguments prove both the support and period
parts of \texttt{hasDepth\_phiA}.

A square flip is precisely forward Fourier exchange, this middle
operation, and inverse Fourier exchange. The other tensor variables
are unchanged. Induction along the actual ordered list
\texttt{flipOrder} therefore preserves both the three clauses of
\texttt{IsTensor} and field membership. At a new vertex its level
is the full integral group, so the required open-subgroup and
integrality hypotheses persist. The original seed has the same
properties because it is the prescribed product of lattice
indicators and determinant characters. This proves arithmetic for
the actual \texttt{Tflip}, not a substitute transform.

Here are the exact field hypotheses in the global application.
The character is
$\nu_{ij,t}=\mu_{0,j}\xi_t/\chi_{m-i}$. If $x=u\pi^a$ then
\[
 |x|^{z/2}=(\sqrt q)^{-az}.                                \tag{TT.7}
\]
Thus \texttt{NumbersIn.chars}, \texttt{sqrt} and \texttt{roots}
supply all coefficients. In particular $q\in\pi\mathcal O_k$,
so $q^e$ kills $\pi^{-e}\mathcal O_k/\mathcal O_k$ and
$\psi(x)^{q^e}=1$ there. This does not identify the rational
integer $q$ with $\pi$, or its valuation with one. The same
calculation, including the normalized half-twist of $\eta$, proves
that the source's \texttt{moved} tensor is $K$-valued after its
specified inverse character factor and frame substitution.

\paragraph{TT5: flat integrals, their phase bounds and their local values.}
For fixed column and row variables $x,y$, the flat integrand is
\[
 e\longmapsto T_g(\operatorname{fill}(x,y,e))
                       \psi(\operatorname{phase}(x,e)),     \tag{TT.8}
\]
where $T_g$ is the actual moved opposite-path tensor. It is locally
constant. It has compact support in the free variables $e$: each
free variable is an entry of a specified \texttt{stepPt} matrix,
and the \texttt{downMap} frame is invertible. If the tensor is
nonzero, its lattice support bounds the entries of every filled
edge after conjugating by fixed matrices. Solving these entry
identities bounds each free variable. A closed bounded box in a
finite-dimensional local vector space is compact. This proves
actual integrability of (TT.8), in contrast to the totalized
multiplicative integral in TT3.

We also record the uniform bound needed for the phase field. For a
fixed moved tensor, all column and row entries visible in its edge
matrices are bounded on the support of the flat integral. The one
remaining last-row direction is bounded by the same translation
argument as TT2: translate the relevant free variables by a small
additive direction. Tensor periodicity leaves the first factor
unchanged, while the second factor is multiplied by
$\psi(a t)$, where $a$ is the last-row entry. A nonzero integral
forces $\psi(a t)=1$ for every $t$ in the chosen small ball.
The annihilator of a local lattice is a lattice; hence $a$ lies in
a fixed bounded ball. This supplies the omitted last-row bound.
Each term of \texttt{phase} is either a matrix entry or a product
of two such bounded entries. The ultrametric triangle inequality
gives, for some integer $N$, the exact implication
\[
 \operatorname{flat}(x,y)\ne0,\quad
 T_g(\operatorname{fill}(x,y,e))\ne0
 \ \Longrightarrow\ \operatorname{phase}(x,e)
                          \in\pi^{-N}\mathcal O_k.          \tag{TT.9}
\]
No bound on $e$ or $x,y$ is asserted when these nonvanishing
conditions fail.

If \texttt{PhasesIn} holds, the integrand is $K$-valued wherever
its tensor factor is nonzero, and it is zero elsewhere. TT1 proves
that every local \texttt{term.datum} value is in $K$. The auxiliary
AA vacuum is a lattice indicator times the standard conductor
phase, and the same compact-box argument proves its integrability.
Its separately prescribed scalar is a product of powers
$q^{n/2}$; it belongs to $K$ by (TT.7) at the AA primes.

\paragraph{TT6: nonzero unfolding terms have finitely many coset labels.}
We give the local bounded-representative argument, since mere local
constancy would not imply finite support on an infinite double
quotient. Write $\operatorname{Bd}(\gamma,M)$ when all entries of
$M$ have valuation at most $\gamma$, and take $\gamma\geq1$.
The elementary block identity behind the argument is this. If
\[
 M=\begin{pmatrix}A&B\\0&D\end{pmatrix},\qquad
 M^{-1}=\begin{pmatrix}A'&B'\\0&D'\end{pmatrix},
\]
and $D,B',D'$ have bound $\gamma$, then
\[
 M=\begin{pmatrix}A&0\\0&I\end{pmatrix}
       \begin{pmatrix}I&-B'D\\0&D\end{pmatrix},\qquad
 \begin{pmatrix}I&-B'D\\0&D\end{pmatrix}^{-1}
       =\begin{pmatrix}I&B'\\0&D'\end{pmatrix}.             \tag{TT.10}
\]
The two right-hand matrices have bound $\gamma^2$. This is exact
block multiplication, using $AB'+BD'=0$ and $D'D=I$.

At one step of the unfolding, perform lattice Gaussian elimination
to write the next matrix as $h=b\kappa$, where $\kappa$ is
integral with integral inverse and $b$ is block upper triangular
for the source's \texttt{midLev}/\texttt{lev} flag. Suppose the
previous matrix and its inverse have a fixed entry bound. A nonzero
flat integral has a point with nonzero integrand, so tensor support
bounds $h^{-1}\operatorname{stepPt}(x,e)$ after multiplying by
the previous matrix. For a free tail direction use
$d=bE_{\mathrm{tail}\,i,\,\mathrm{tail}\,(i+1)}t$, with
zero-based tail indices. This $E$ is rectangular: its tail has
$c$ rows and $c+1$ columns, exactly as in \texttt{stepPt}.
Thus for the last row $i=c-1$ its column $i+1=c$ still exists.
For $c=1$ the same translation uses row zero and column one,
and bounds the sole tail diagonal; no endpoint is omitted.
Write $\rho=\lfloor c/2\rfloor$ and $p=\rho+1$. The exact free pattern is: no core--core entries;
tail row $l$ and a core column when $l+1<p$; a core row and tail
column $j$ when $p<j+1$; and tail row $l$, tail column $j$ when
$l<j$. The auxiliary translation uses only the free entries,
not the extra \texttt{Kept} entries. Its sole possible nonzero
column is tail $i+1$. Upper triangularity for
$\operatorname{midLev}(\rho)$, which assigns level $2\rho+1$
to the core and $2(l+1)$ to tail $l$, forces a nonzero tail row
to have $l\leq i$, and a nonzero core row to have $i\geq\rho$.
These are exactly the required free inequalities in column $i+1$.
All core columns are zero. Moreover,
\[
 h^{-1}d=\kappa^{-1}E_{\mathrm{tail}\,i,\,\mathrm{tail}\,(i+1)}t.
\]
Since $\kappa^{-1}$ is integral, a fixed ball for $t$ makes this
a tensor period after multiplication by the bounded previous
matrix. The ball is independent of the possibly unbounded $b$.
The source's character reads the sum of entries
$(\mathrm{tail}\,l,\mathrm{tail}\,(l+1))$, so on $d$ it is
exactly $b_{ii}t$. Zero core columns keep the additional row
terms in the phase unchanged. Translating the integral then forces
\[
                    \psi(b_{ii}t)=1                         \tag{TT.11}
\]
throughout that ball. The annihilator bound gives a common bound
for every tail diagonal entry of $b$. The fixed entries of
\texttt{stepPt} likewise bound the tail columns of the inverse
of the diagonal block part of $b$. Removing its flag-unipotent
part and applying (TT.10) yields
\[
                  h=r w,\qquad r\in R_j,\qquad
     \operatorname{Bd}(C,w),\quad\operatorname{Bd}(C,w^{-1}), \tag{TT.12}
\]
where $R_j$ is exactly the core general-linear group times the
specified symmetric unipotent group. No enlargement to all
unipotent matrices has been made. This argument also covers an
empty core; a zero tail has no new tail equations.

The set of $w$ in (TT.12) is compact. To see closedness without
assuming inversion extends to singular matrices, view it as the
projection of the compact set of pairs $(w,v)$ in the two bounded
matrix boxes satisfying $wv=vw=I$. If $U$ is the prescribed open
right level, its image in $R_j\backslash\mathrm{GL}/U$ is finite:
the open right cosets cover that compact set and a finite subcover
suffices. Induct on the vertices. There is only the fixed initial
label; finitely many previous representatives have one common
entry bound; (TT.12) then leaves finitely many next labels. The
finite product of these finite label sets contains every nonzero
unfolding term. Taking the product over the finitely many active
places proves the literal \texttt{S.supp} is finite.

\paragraph{TT7: pure-coset expansion and fixed-depth uniqueness.}
For a fixed nonzero unfolding label, TT5 gives bounded support and
a sufficiently small common translation period for its flat datum.
After clearing the finitely many frame and half-lattice denominators,
choose one positive integer $k$ so that each relevant half-point
set is partitioned into finitely many cosets
\[
                  x_0+k\Lambda_{\rm half},\qquad
                       kx_0\in\Lambda.                    \tag{TT.13}
\]
Here the source's idempotent is fixed, and
$\Lambda_{\rm half}$ means its \texttt{IsHalf} subgroup; it is
not an arbitrary choice of a half of the ambient vector space.
Existence of the finite representatives is the finite quotient of
a bounded fractional lattice by a full integral sublattice. CRT
and localization identify this quotient with the actual local
cosets in the given frames. There is no uniform bound for $k$
claimed in this paragraph.

For each tuple of representatives take the product of its top and
row coset indicators and assign to it the value of the flat datum
at the tuple of centers. Only finitely many coefficients are
nonzero. Periodicity makes this a pointwise expansion on half
points; off the support both sides vanish. With the prescribed
prime set, idempotent and $z=1$, these are actual
\texttt{FinDatum}s and give \texttt{IsCosetExpansion}. In
particular some positive depth exists, so the source's least depth
\texttt{expDepth} exists and is positive.

The following stronger uniqueness is useful. Two depth-$k$ coset
indicators with the same idempotent that are nonzero at a common
point have the same coset. With the same prime set and $z=1$ they
are the same \texttt{FinDatum}. Thus the product indicators of
different pure data have disjoint supports. Evaluating any
expansion at a tuple of centers of a nonzero term $d$ gives
\[
       \tau(d)=\prod_w
        \operatorname{datum}_w(x_w,y_w).                   \tag{TT.14}
\]
It also forces that same $d$ to occur in every other expansion at
the same depth, with the same coefficient. Consequently the union
of the supports of \emph{all} depth-$k$ expansions is contained
in one chosen finite support. This proves the literal
\texttt{finite\_pureDat\_of\_depth}; it does not use the false
inference that a union of arbitrary finite sets is finite.

\paragraph{TT8: expansion independence and the actual SumOK supplier.}
The class pairing is independent even of expansions not restricted
to the special disjoint basis (TT.13). At fixed idempotent, each
finite theta-vector construction is linear in its finite datum
function: apply the prescribed finite Heisenberg sum to that
function. The theta-vector characterization and its normalization
identify this explicit vector with the chosen $Z.\mathrm{vec}$.
Thus the kernel operator is linear in each datum function; after
composition down the top and row chains, the final class pairing
is multilinear in the $2m$ datum slots.

To justify the tensor argument without an infinite-dimensional
assumption, collect the finitely many functions occurring in the
two expansions in each slot and take their finite-dimensional
spans. Point evaluations separate each such span: if every
function in a subspace vanished at every point, that subspace
would be zero. Choose finitely many evaluations giving an
injective coordinate map by successively reducing the kernel.
Their tensor product is injective. Therefore a finite sum of
pure tensors whose associated product function vanishes pointwise
is the zero tensor. Every multilinear functional annihilates it.
Apply this to the difference of the two expansions. Both vanish
off the specified half points, so their equality on half points
is sufficient. Summing over classes preserves the equality.

The remaining class conditions are actual consequences of the
geometric setting. Every top and row space used here is positive
definite at all archimedean places. Its level is open compact:
finite local projected levels are open and bounded, and the
outside lattice stabilizers are compact. The lattice argument of
CR3 gives finitely many classes; the automorphism group of a
positive definite integral lattice is finite, so every stabilizer
order is a positive integer. Accordingly
\[
             \operatorname{invGam}(c)
                    =\prod_j|\Gamma_{c,j}|^{-1}\in\mathbb Q^\times.
                                                                  \tag{TT.15}
\]
Trivial stabilizers are not required. The displayed weight lists
in the grid give the indicated interlacing and common determinant
shift by substitution. Reading the row coefficient space uses its
actual \texttt{funEquiv}; it is an isomorphism of the stated
polynomial representations, not an assertion about equality of
two integral lattices.

TT4--TT7, the preceding multilinearity, and these class facts
supply, respectively, all eleven fields of \texttt{SumOK}:
\texttt{flat}, \texttt{flatA}, \texttt{fin}, \texttt{ex},
\texttt{vec}, \texttt{indep}, \texttt{cls}, \texttt{gam},
\texttt{wts}, \texttt{br}, and \texttt{read}. They hold for every
literal choice at the given Hida setting. Thus \texttt{SumOKK}
holds with both its additional largeness thresholds equal to zero.
The source's \texttt{rowSum\_of\_hypotheses} has a nominal
\texttt{FunPr.Hypotheses2} argument, but its proof does not use it;
the proof above likewise uses none of that additional package.

\paragraph{TT9: rational polynomial coordinates, operators and branching.}
Let $b_1,\ldots,b_r$ be the prescribed bideterminant basis of a
coefficient space $VV_\lambda$. Every $b_i$ has integer
coefficients. A polynomial in its complex span has all coefficients
in $K$ if and only if its $b_i$-coordinates lie in $K$. For the
nontrivial implication choose finitely many monomials such that
the corresponding coefficient matrix of the $b_i$ has a nonzero
$r\times r$ minor. That minor has rational entries and inverse
over $\mathbb Q$. Reading those polynomial coefficients recovers
all $b_i$-coordinates in $K$. This is finite-dimensional field
descent, with no assertion about unimodularity of that minor.

A rational frame acts by linear substitution in the variables and
an integral power of a nonzero determinant. All its coefficients
lie in $K$ after the fixed embedding of $E$; inverse transpose
uses the inverse of that same frame. The matrices $Z.\mathrm{kk}$
are in $K$ under \texttt{JetsIn}. Hence the exact top and row
chains preserve $K$-polynomial coefficients. The row chain starts
with the source's \texttt{funOf} and the same inverse rational
frame, in that order.

The final branching pairing also preserves $K$. The Fock pairing
has rational diagonal coefficients (products of factorials) in
the monomial basis. Restriction, adjugate substitution and
bideterminant operations have rational coefficients. Orthogonal
projection onto the indicated rational polynomial subspace is
computed by inverting its rational nondegenerate Fock Gram
matrix, and is therefore rational. The maps \texttt{covOf} and
\texttt{funOf} are inverse rational maps on their stated spaces.
Finally the extremal normalization \texttt{extNormB} is the
nonzero rational Fock self-pairing of the fixed extremal
polynomial. Dividing by it proves
\[
       U,V\hbox{ with coefficients in }K
       \quad\Longrightarrow\quad\operatorname{brPair}(U,V)\in K.
                                                                  \tag{TT.16}
\]
These statements concern the bilinear Fock/branching conventions
in the source, including their specified conjugate identification;
no field is silently assumed stable under arbitrary complex
conjugation. All maps used in (TT.16) are the rational polynomial
maps after that identification.

For the actual source vector $x$ in its rational form, VE5--VE11
and the given model identity express its value at a finite point
as a rational coefficient vector multiplied by the actual level
character. That character is a product of the type characters and
inverse twist characters; the initial flag character is one.
Their values lie in $K$. Applying the fixed rational coefficient
identification gives $\texttt{formC}(j_X x)$ with polynomial
coefficients in $K$. This checks the same-object input to the row
chain; a rank equality alone would not suffice.

\paragraph{TT10: the actual global sum is rational over the given field.}
Choose a depth-$\mathrm{expDepth}$ expansion $\tau_\beta$ for
each nonzero unfolding label. TT8 permits this choice inside the
literal definition of $v_K$. That definition is
\[
\begin{split}
 v_K(t,x,a)={}&s_{AA}
 \sum_{\beta\in\operatorname{supp}}
       \left(\prod_w c_w(\beta_w)\right)
 \sum_{c\in\operatorname{Classes}_\beta}
       \operatorname{invGam}(c)\\[-2mm]
 &\hspace{6mm}\cdot\sum_{d\in\operatorname{supp}\tau_\beta}
       \tau_\beta(d)\,
       \operatorname{brPair}\bigl(
        \operatorname{colMap}_d\Psi_a(c_{\rm top}),
        \operatorname{rowMap}_d\operatorname{formC}(j_Xx)(c_{\rm row})
                              \bigr).
\end{split}\tag{TT.17}\]
Every displayed sum is finite by TT6--TT8. The local coefficients
$c_w(\beta_w)$ are the prescribed rational quotient volumes.
TT5 gives $s_{AA}\in K$, and (TT.14) and TT5 give
$\tau_\beta(d)\in K$. Equation (TT.15) treats the stabilizers.
The test column $\Psi_a$ is the scalar extension of its actual
coefficient-field vector, so TT9 applies on both sides of the
pairing. Every term of (TT.17) is consequently in $K$.

This proves the literal \texttt{SumRatK}, for the original choice
and the original reference field, under its stated
\texttt{FieldOK}. It uses TT8 for \texttt{SumOKK}. No new
largeness threshold is needed, and no global period-integral
convergence or identity with (TT.17) is inferred from this
arithmetic calculation.

\paragraph{TT11: the character, root and phase portions of NumbersIn.}
There is one number field containing all local values of each
algebraic Hecke character used here. A direct argument avoids
assuming that infinitely many algebraic numbers generate a finite
extension. Choose an open unit subgroup on which the finite
character is trivial, and let $\mathfrak f$ be its modulus.
The relevant ray class group is finite. For a finite idele, divide
by a chosen representative of its ray class and a principal idele;
the remaining unit lies in that open subgroup. The principal
value is the inverse infinity-type monomial, hence belongs to the
normal closure of the number field defining the embeddings. The
finitely many chosen ray-class values are algebraic: raise each
to the order of its ray class, and use that same principal-value
formula. Adjoin these finitely many roots. This proves the claimed
single field. Finite-order characters are simpler: all their
values belong to one cyclotomic field.

Apply this to $\mu_0$, the grid characters $\chi_i$,
$\chi_\psi$, and the finite-order twists and discriminant
characters. The actual integral infinity exponents follow from
\texttt{LevCh.Is}; the half-twist of $\eta$ is kept in its
\texttt{ValuesIn}(1/2) convention, supplied by the same-HFrame construction HF16 and TL14. We do not adjoin $\sqrt q$ for infinitely
many rational primes. Instead, \texttt{NumbersIn.sqrt} concerns
only the finite union of AA and active local primes. Adjoin those
finitely many positive square roots, the eighth roots of unity,
and the roots of order $q_w^{e_w}$ at the active places. This
constructs the literal \texttt{CharsIn} and \texttt{RootCl}.

For phases, (TT.9) and $q_w\in\pi_w\mathcal O_w$ give
$\psi_w(z)^{q_w^{N_w}}=1$ on the required support. There are
finitely many twists, finite grid data, unfolding labels and active
places, by the actual finite quotient defining \texttt{G.Datum}
and TT6. Take the product of their positive orders and adjoin
its roots of unity. This proves \texttt{PhaseCl}, including its
two nonvanishing conditions. The fields just constructed depend
on the fixed actual data; their compositum is finite over
$\mathbb Q$. These are the actual \texttt{NFChars},
\texttt{NFRoot} and \texttt{NFPhase} suppliers, with zero
additional largeness thresholds.

\paragraph{TT12: the moved lattices and vacuum use the same rational frame.}
At a definite node let $B$ be the normal datum's rational frame,
$B^*JB=I$, let $u$ satisfy the actual lattice identity
$\widehat\Lambda=uB\widehat{\mathcal O}_E^n$, and let $x$
be the finite point. In $B$-coordinates the
actual moved lattice is
\[
 L_x=E^n\cap a\widehat{\mathcal O}_E^{\,n},\qquad
             a=B^{-1}(x u)B.                               \tag{TT.18}
\]
The order $xu$, and the inverse on the left of $B$, are forced by
$Bv\in\operatorname{globLat}(x,\Lambda)$: its defining condition
is $x^{-1}Bv\in\widehat\Lambda$, and
$\widehat\Lambda=uB\widehat{\mathcal O}_E^n$.
Both $x$ and $u$ preserve $J$, so direct multiplication gives
$a^*a=I$. Localization and the paired-prime duality of NB1 show
that each localization of (TT.18) is self-dual. The finite
intersection/localization argument LP2 then gives a full global
self-dual lattice. This proves the actual
\texttt{srcLat}/\texttt{tgtLat} statements with the source's
swap and sign at a row edge; it does not choose new frames.

On the completed Hom lattice, the action is
$X\mapsto g_T Xg_S^{-1}$, with the source's transpose convention
when written as row coordinates. Substitution of the two lattice
conditions proves inclusion in the new Hom lattice. Apply the
inverse action for the reverse inclusion. Consequently the actual
Weil image of the old vacuum has the new lattice translation and
half-translation laws, by Heisenberg covariance. At infinity the
acting point is the identity, so its annihilation operators commute
with this action; the Fock vacuum equations persist. Its norm is
one because the supplied Weil action is unitary for the same
fundamental-domain measure. These prove the full moved-vacuum
conditions, including normalization, from the given Weil datum.
They do not yet prove that the scalar comparing two normalized
vacuum vectors is algebraic.

\paragraph{TT13: normalized jets from finite arithmetic on the same vacuum.}
We isolate a narrower remaining arithmetic assertion than
\texttt{NFJet}. For every actual definite edge $p$ and finite
pair $g$ used in the class sums, require
\[
 W_p(g)\operatorname{vac}_p
       \in\operatorname{span}_{\overline{\mathbb Q}}
       \{\operatorname{Heis}(v)\operatorname{vac}_p:
                                      v_\infty=0\}.         \tag{TT.19}
\]
The span means a finite linear combination. It concerns this
particular $W_p$, its given splitting characters, and its given
unit vacuum. TT12 by itself does not imply (TT.19): an unspecified
complex unit scalar on a line need not be algebraic. We retain
(TT.19) exactly where it is used below.

The normalized base entire function $e_{0,p}$ in the normal datum
has the source's \texttt{ThetaVectorUpToUnits} description: an
algebraic unit times a finite positive intertwiner applied to the
product one-variable theta function in a rational orthonormal
basis. TA10, VI4 and VI9--VI10 imply that every rational
Heisenberg translate of this function has algebraic Taylor
coefficients. This is the direct finite-translate argument. It
does not require the stronger general assertion that an arbitrary
algebraic-coefficient entire function remains so after every
translation.

For clarity, finite-adelic translations have exactly the required
rational-coordinate interpretation. If $v_\infty=0$, choose
$n>0$ with $nv$ in the completed Hom lattice. Additive approximation
and CRT choose a rational vector in the prescribed finite coset.
Subtract it and a lattice vector from $v$; what remains is its
opposite rational archimedean coordinate. Heisenberg composition
and the vacuum's half-lattice law introduce values of the standard
character on rational finite adeles. Multiplication by an integer
puts each argument in the character's conductor lattice, so these
values are roots of unity. The exact Fock derivative formula is
therefore a root of unity times the derivative of one rational
translate of $e_{0,p}$, multiplied by the original period factors.
This also proves the source's finite-translation jet statement
without changing the normalized half-point or $\delta$.

A coset datum with $z=1$ is an actual finite algebraic Heisenberg
combination of the vacuum. To see this, take a common finite
lattice quotient on which its indicator and character are
constant, and apply the finite Fourier identity (TT.6). The
coefficients are rational multiples of roots of unity; the
standard theta-vector characterization identifies this finite sum
with $Z.\mathrm{vec}(d)$. Heisenberg covariance and (TT.19) give
the same conclusion after the actual finite point acts on that
vector. It follows that its raw jet table has the form
\[
 J(m)=c_{0,p}\left(\prod_\tau
        \operatorname{opp}_p(\operatorname{period}(\delta,\tau))
                      ^{d_\tau(m)}\right)A(m),
                   \qquad A(m)\in\overline{\mathbb Q}.     \tag{TT.20}
\]
Here $\operatorname{opp}_p$ is precisely identity or conjugation
for that oriented edge, and $d_\tau(m)$ is the total degree at
$\tau$. Both preserve algebraicity. No arbitrary field embedding
is interchanged with a limit, derivative or integral.

The jet operator is homogeneous of the specified Fock degree
$f_{p,\tau}$ at each place. Thus (TT.20) factors out
$c_{0,p}\prod_\tau\operatorname{opp}_p
(\operatorname{period}(\delta,\tau))^{f_{p,\tau}}$.
The actual scalar $Z.\mathrm{scal}(1,p)$ is its inverse.
All periods and $c_{0,p}$ are nonzero in the normal datum. The
factors cancel exactly. The remaining operator is obtained from
$A(m)$ by rational Fock differentiation, rational projection and
bideterminant coordinates (TT9). Hence every actual normalized
entry $Z.\mathrm{kk}(d,x,y)_{ab}$ is algebraic. Alternatively,
the literal \texttt{DictAt} formula supplies (TT.20) directly:
its $\kappa$, $r_i$ and rational translates are algebraic, and
its same $c_0,\delta$ cancel in precisely this calculation. The
extra norm, denominator and all-embedding estimates in
\texttt{DictAt} are not needed for this qualitative conclusion.

\paragraph{TT14: finitely many jet entries give one field.}
For the fixed choice, consider all twists, all finite grid data,
all nonzero unfolding labels, all classes, all top and row edges,
and all their bideterminant matrix indices. These are finite by
the actual finite datum quotient and TT6--TT8. For each label use
the least depth \texttt{expDepth}. TT7 shows that the pure data
appearing in \emph{any} expansion at that depth form a finite
set. Thus the set $E_{\rm jet}$ of all entries required by
\texttt{JetsIn} is finite, despite its universal quantifier over
expansions. Under (TT.19), or the stated qualitative dictionary,
TT13 shows every member of $E_{\rm jet}$ is algebraic. Then
\[
                  M_{\rm jet}=\mathbb Q(E_{\rm jet})        \tag{TT.21}
\]
is a number field and proves the literal \texttt{JetCl}. In
particular a same-data supplier of (TT.19) would close
\texttt{NFJet}; no uniform degree bound is claimed here. This
argument supplies the finite-generation step that pointwise
algebraicity alone does not supply.

\paragraph{TT15: algebraic compact averaging and one surviving degree.}
Let $C$ be a compact group with probability Haar measure, and
$\rho:C\to\mathrm{GL}(V)$ a continuous finite-dimensional
complex representation. Suppose a dense subset of $C$ has
matrices over $\overline{\mathbb Q}$ in the chosen basis. Then
\[
                         P=\int_C\rho(g)\,dg               \tag{TT.22}
\]
has algebraic entries. Here is an algebraic proof of this fact,
which does not presume that integrals of algebraic-valued
functions are algebraic. Put $F=\overline{\mathbb Q}$, and form
over $F$ the span $D_F$ of all columns of $\rho(\gamma)-I$ and
the common kernel $I_F$ of those matrices, for the dense set.
Finite-dimensionality reduces both constructions to finitely many
matrices. Their complex scalar extensions are respectively
$D=\operatorname{span}_{\mathbb C}(\rho(g)-I)V$ and $V^C$:
for $D$, its finite-dimensional span is closed, so density puts
all $g$ in it; for $V^C$, use continuity of $g\mapsto\rho(g)v$.
Haar invariance shows that $P$ is the identity on $V^C$ and zero
on $D$. Also $v-Pv$ lies in $D$, by integrating in this closed
finite-dimensional subspace. Therefore
$V=V^C\oplus D$. The same direct-sum decomposition holds over
$F$, by ranks of finite matrices. Its projection has coefficients
in $F$, proving (TT.22). The zero-dimensional case is included.

For a supplied actual bottom triangle, the real group on the
source side is a product of compact unitary groups. In its
standard frames the Fock vacuum character is
\[
            \chi(g)=\prod_\tau\det(g_\tau)^{\,N-k_\tau},
       \qquad b_\tau=k_\tau-N/2,\quad k_\tau\in\mathbb Z,
                                                                  \tag{TT.23}
\]
where $N$ is the target rank and $b_\tau$ is the infinity
exponent of the specified splitting-character ratio. The parity
restriction of that ratio gives the displayed integrality.
The sign is the negative source-space sign: the positive
metaplectic half-trace and the splitting-character exponent
combine to $N/2-b_\tau=N-k_\tau$. More explicitly, diagonalize a
unitary source matrix with eigenvalues $e^{\mathrm i\theta_j}$.
For one oscillator,
$H=(-d^2/dx^2+x^2)/2$ satisfies
$H(e^{-x^2/2})=\tfrac12e^{-x^2/2}$, since the second derivative
is $(x^2-1)e^{-x^2/2}$. In this negative-source polarization each
of the $N$ Gaussian coordinates in source direction $j$ therefore
contributes the rotation factor $e^{\mathrm i\theta_j/2}$; the
prescribed splitting character contributes
$e^{-\mathrm i b_\tau\theta_j}$.
Their product is $e^{\mathrm i(N/2-b_\tau)\sum_j\theta_j}$.
The vacuum is unchanged by a unitary change of the diagonalizing
basis, and a creation monomial is carried to the corresponding
linear substitution. The integer exponent makes the result
independent of the angles chosen. This gives (TT.23) on the
whole compact unitary group, including its special-unitary part. A
finite-order twist is trivial on its connected real subgroup, so
the same $\chi$ works for every twist $t$.

On monomials with place-degree vector $d$, the resulting finite
matrix representation is $\chi(g)\operatorname{subst}_d(g)$.
Algebraic-entry unitary matrices are dense: use the Cayley
transform on skew-Hermitian matrices with algebraic entries,
which is dense in the open chart without eigenvalue $-1$.
Its matrices are algebraic, including negative determinant
powers in (TT.23). Thus TT15's first argument applies to every
degree block. At a central scalar $z$ in one real place, the
representation is multiplication by $\chi(z)z^{d_\tau}$.
Choose $z=e^{\mathrm i}$, which has infinite order. Haar
averaging vanishes unless
$\chi(z)z^{d_\tau}=1$ at every place. There is at most one
nonnegative degree vector $d_0$ satisfying these equations. If
there is none, all averages are zero and any fixed $d_0$ may be
used with zero coefficients. Otherwise every averaged kernel is
an algebraic linear combination of the finitely many kernels of
degree $d_0$. This is the actual fixed-degree conclusion used in
\texttt{avg\_kernel\_alg}, on the same vacuum and the same
probability Haar measure.

\paragraph{TT16: normalized bottom periods and their finite field.}
Fix a literal \texttt{IsBottomTri} with its actual Haar measures
and finite-order twists. Retain the following finite-Weil input:
for both source sides, its standard theta datum and its images
under finite adelic pairs are finite algebraic combinations of
finite Heisenberg translates of the corresponding fixed vacuum,
as in (TT.19). This is the content whose source producer is
\texttt{moved\_translate\_algComb}; its existence is not a
consequence of compact averaging alone.

For a fixed place degree $d$, VI9--VI11 and the finite-translation
argument of TT13 give a common nonzero period
\[
       \Omega_d=\overline{c_0\prod_\tau
                      \operatorname{period}(\delta,\tau)^{d_\tau}}
                                                                  \tag{TT.24}
\]
for the negatively oriented source; use the corresponding
unconjugated expression on the other source side. The unspecified
unit relating the actual normalized vacuum to its entire model
is included in $c_0$, so it is not asserted algebraic. Nonzero
vacuum norm and nonzero periods imply $\Omega_d\ne0$.
The retained finite-Weil assertion and the explicit finite
Fourier expansion of the standard theta datum show that every
finite-point monomial kernel of degree $d$, divided by
$\Omega_d$, is algebraic.

There is a common compact open finite level for these finitely
many datum functions and twists. Lattice stabilizers preserve the
vacuum line, and their action on it is a continuous character.
A continuous character of a profinite group into the circle has
finite image: its compact image is a closed subgroup of the
circle, and a quotient of a profinite group is totally
disconnected; the only such closed subgroups of the circle are
finite. Shrink the congruence subgroup to kill this character,
and to preserve the finitely many Heisenberg translations and
their conductor phases. This supplies a level fixing all the
finite datum functions. Polynomial creation operators commute
with finite-place action, so the level does not depend on the
monomial.

Let $z_l$ be the finite set of definite adelic class
representatives, let $w_l>0$ be their finite stabilizer orders,
and put $v=\mu(J)>0$, where $J$ is the product of that finite
level and the full compact real group. With probability Haar
measure $dg$ on the latter, the exact quotient formula is
\[
 I_t(P,h)=v\sum_l\frac1{w_l}
                    \int_C\operatorname{ker}_t(P,h)(z_lg)\,dg.
                                                                  \tag{TT.25}
\]
It follows by decomposing the adelic quotient into these classes;
the finite stabilizer makes each fiber a quotient of multiplicity
$w_l$, giving its reciprocal. Compactness gives convergence for
each continuous kernel. The number $v$ is the actual measure
normalization and need not be rational. TT15 reduces the integral
to the single degree $d_0$ with algebraic coefficients. Thus
\[
                         I_t(P,h)/(v\Omega_{d_0})
                                      \in\overline{\mathbb Q}     \tag{TT.26}
\]
for monomials $P$ and finite points $h$. This produces one
nonzero period for all twists and all degrees, including the
case where every period integral is zero.

We must still pass from algebraicity to a number field. Left
rational invariance, the common finite level, and the finite
definite class set reduce $h$ to finitely many representatives,
up to an archimedean point inverse to the archimedean component
of a rational unitary element. Its matrix coefficients lie in
the normal closure of $E$. Fock substitution and the integer
determinant powers in (TT.23) therefore have coefficients in one
fixed number field. Only finitely many monomials can have a
nonzero integral: central real scalars force the degree vector
of TT15, or the same argument applied before averaging for each
of the finitely many twists. Consequently all integrals are
linear combinations, over a fixed number field $M_0$, of finitely
many values $I_t(P_b,h_a)$. Adjoin their normalized algebraic
values from (TT.26). The resulting finite extension contains
all $I_t(P,h)/(v\Omega_{d_0})$. Repeat independently for $I'_t$;
its period and its Haar volume may differ. The compositum proves
the exact \texttt{SumsIn} assertion.

This proves the substantive definite-period implication from the
specified finite-Weil arithmetic. To obtain \texttt{NFBot} for
all squares, one also needs an actual bottom triangle at each
square. The source builds standard positive spaces of ranks
\texttt{botBig}, \texttt{botSmall}, and their sum, then chooses
Weil representations with its displayed parity-twisted characters.
The lattice and rank construction is elementary; the existence
of those additional full adelic Weil representations is retained
here. Their existence is not supplied merely by the grid's
already chosen $W$.

\paragraph{TT17: common rational periods force the bottom scalar ratio.}
We prove the precise algebraic comparison used after TT16. Suppose
there is one complex vector space $\Pi$, complex-linear maps
$p:P\twoheadrightarrow\Pi$, $p':P'\to\Pi$ with nonzero image,
and injective maps $J_t:\Pi\to\{\text{functions of }g\}$ such
that
\[
 I_t(P)(g)=\lambda_tJ_t(pP)(g),\qquad
 I'_t(P')(g)=J_t(p'P')(g),\qquad\lambda_t\ne0.              \tag{TT.27}
\]
These are the substantive factorization hypotheses; they are
kept on the same bottom triangle. Assume its monomial periods
are rational over $K$ after division by nonzero $\Omega,\Omega'$,
and finite-point vanishing of a difference in (TT.27) implies
vanishing at every point. Then
\[
                       \lambda_t/\lambda_{t_0}\in K^\times.\tag{TT.28}
\]

Choose $P'$ with $p'P'\ne0$ and $P$ with $pP=p'P'$. There is a
finite point $g_0$ where $I'_{t_0}(P')(g_0)\ne0$, by injectivity
and the finite-point determination hypothesis. Choose a
$K$-linear functional $\varphi:\mathbb C\to K$ nonzero on this
value, by extending a basis of its one-dimensional $K$-span.
No continuity of $\varphi$ is assumed or needed. For a polynomial
$Q=\sum a_mX^m$, put
$\operatorname{rat}_{\varphi,w}(Q)=\sum\varphi(wa_m)X^m$,
viewing its coefficients in $\mathbb C$ through $K\subset\mathbb C$.
If a complex-linear functional $\ell$ has
$\ell(X^m)\in\Omega K$ for every $m$, direct finite expansion gives
\[
 \ell(\operatorname{rat}_{\varphi,w\Omega}(Q))
                         =\Omega\varphi(w\ell(Q)).          \tag{TT.29}
\]
Set
$Q=\operatorname{rat}_{\varphi,\lambda_{t_0}^{-1}\Omega}(P)$
and $Q'=\operatorname{rat}_{\varphi,\Omega'}(P')$.
At $t_0$, (TT.29), finite-point determination and injectivity yield
\[
             \Omega'\lambda_{t_0}pQ=\Omega p'Q',
                              \qquad p'Q'\ne0.              \tag{TT.30}
\]
The second assertion follows by evaluating at $g_0$ and using
$\varphi(I'_{t_0}(P')(g_0))\ne0$.
For each $t$ choose a finite point with $I'_t(Q')\ne0$ and write
$I_t(Q)=\Omega z$, $I'_t(Q')=\Omega'z'$ with $z,z'\in K$,
$z'\ne0$. Apply $J_t$ to (TT.30) and multiply by $\lambda_t$.
Cancelling the nonzero periods gives
$\lambda_{t_0}z=\lambda_tz'$, hence (TT.28).

For the actual definite continuous automorphic kernels,
finite-point determination has a direct justification. Rational
unitary points are dense in the compact real group: the Cayley
chart is rational on the vector space of skew-Hermitian matrices,
and weak approximation approximates their entries at all real
places. For a fixed finite component $g_f$ and rational points
$\gamma_n$ with $\gamma_{n,\infty}\to g_\infty$, left
invariance gives
$F(\gamma_{n,\infty},g_f)=F(1,\gamma_{n,f}^{-1}g_f)=0$.
Continuity gives $F(g_\infty,g_f)=0$. No convergence of the finite
components is required. Compactness of the definite quotient
also gives the integrability needed to use the monomial period
identities. Thus TT17 proves the actual
\texttt{Factor.ratIn}/\texttt{ratIn\_of\_definite} implication;
it does not construct the common factorization $\Pi,p,p',J_t$.

\paragraph{TT18: the exact arithmetic consequence of a matching tower.}
Fix a square $(i,j)$ and retain its actual nonzero square
constants $c_t$, bottom factorization (TT.27), and the source's
tower identity. Write $p=\texttt{pp}(j)$,
$q=\texttt{qq}(j)$, $e=\texttt{eTop}(i,j)$, and let
$\Lambda_t$ be the analytic continuations of the finite Hecke Euler products
for the actual characters. The identity is
\[
 \frac{c_t}{c_{t_0}}\prod_{a=0}^{q-1}
       \frac{\Lambda_t(e+a+1)}{\Lambda_{t_0}(e+a+1)}
 =\mathrm{embLL}(\zeta)\left(\prod_{b<p+q}u_b\right)
       \frac{\lambda_t}{\lambda_{t_0}}
       \prod_{a=0}^{p-1}
       \frac{\Lambda_t(e+q-a)}{\Lambda_{t_0}(e+q-a)}.        \tag{TT.31}
\]
Every $u_b$ is the literal \texttt{StepRatio} with final argument
one, and hence lies in $K$ by its displayed coefficient-field
factorization. The away-unit $\zeta$ is an element of the same
coefficient field. Under \texttt{NumbersIn.ratios}, all displayed
$L$-ratios lie in $K$. The exact tower-number inequalities put
their arguments between $1$ and \texttt{halfRank}, and hence
within the upper bound \texttt{bigN}+1 required by that field
clause. The supplied \texttt{LNe} is stronger: it gives
nonvanishing at every integer from $1-\texttt{e0}$ through
\texttt{halfRank}. It therefore makes every denominator and
every factor on the left of (TT.31) nonzero.

Use the bottom triangle selected by \texttt{NumbersIn.bottom},
not an unrelated one. The separate existence input supplies a
witness that this selection is allowed; its factorization is
assumed for every valid bottom triangle. TT17 then gives
$\lambda_t/\lambda_{t_0}\in K^\times$ for the selected one.
Divide (TT.31) by its nonzero left product. This proves the actual
$\texttt{sqRatio}(i,j,t)\in K$. Taking the maxima of the original
constant, tower and nonvanishing thresholds gives the literal
\texttt{RatioK} implication, with all original inputs retained.
In particular a merely algebraic numerator is insufficient, and
qualitative nonvanishing at $s=1$ for an unrelated finite-order
character does not supply this critical interval.

\paragraph{TT19: a finite coefficient field and its same-choice boundary.}
For a fixed actual choice, TT11 supplies its character, root,
square-root and phase fields. Under the exact finite-Weil input
(TT.19), TT14 supplies its jet field. Under the bottom existence
and finite-Weil inputs of TT16, it supplies its bottom period
fields. We can supply the ratio field of \texttt{LRatiosIn} on
this same grid. The characterization \texttt{IsHeckeL} fixes
the finite Euler products on $\Re(s)>1$, so analytic continuation
identifies any two witnesses. There are finitely many twists,
squares and integer arguments $1\le s\le\texttt{bigN}+1$.
For each square, its characters are finite-order twists of the
reference character with the same positive angular weights $2m_\tau$.
The grid gives $m_\tau>s$ on this range, so both gamma arguments
$s+m_\tau$ and $1-s+m_\tau$ are positive. Blasius's critical-value
theorem applies to the finite functions at these integers.
The Euler product for $s>1$ and Hecke nonvanishing at one give
nonzero values. In the quotient of the two critical-value formulas
the common Tate factor cancels. By multiplicativity of CM periods,
the period quotient is the period of a finite-order character up to
an algebraic factor. If the character has order $n$, this period
raised to the $n$th power is algebraic. Thus the critical quotients
are algebraic, and adjoining them gives a number field.

The compositum of these finitely many number fields is finite and
satisfies all eight fields of \texttt{NumbersIn}: characters,
eighth roots, local depth roots, phases, the specified square
roots, jets, critical ratios, and bottom sums. Field inclusion
preserves each clause. For the existential bottom clause use the
same chosen bottom triangle and the same two periods after
inclusion. For jets use the same $Z.\mathrm{kk}$ entries and
least-depth expansions. For phases use the same local integral,
frames and representatives. This proves the actual finite-field
join, rather than constructing eight unrelated choices.

To conclude literal \texttt{FieldOK}, this compositum must be
embedded in the range of the \emph{same} \texttt{embLL} used by
$R$ and $\mathcal C$. If that embedding is supplied, all eight
clauses transfer immediately and TT10 applies. If the coefficient
field is to be enlarged, the source's \texttt{LAgree},
\texttt{GAgree}, \texttt{WAgree}, \texttt{SAgree} and flip
agreements must be provided. They preserve the preliminary datum,
the actual HFrame and reference twist, every space/lattice and
frame, the character functions, theta realization, finite point
and jet data, and the actual local Fourier integrals. Merely
knowing that a finite extension exists does not prove these
same-object agreements or construct \texttt{Choice.Is}; in
particular \texttt{RealOK} is the additional separation condition
in that record. The producers \texttt{exField\_of\_keyFin\_exKey}
and \texttt{dat\_choice\_exists} remain conditional on these
compatible supplies.

On the indicated Hida grids, the increasing-rank first-occurrence
argument gives convergence of the actual iterated theta/Bessel tests.
For \texttt{InjK} one still needs the finite catalogue to detect
every nonzero input. In \texttt{Formula2K}, the good-input
projection, the definite period-to-jet evaluation and the mixed
cell measure factors have been identified. The square matching
and the remaining mixed unfolding concern the prescribed tests.
For \texttt{RatioK} the bottom factorization and the two
character-uniform tower identities still require their analytic
comparisons. Critical nonvanishing follows from Hecke's theorem
and the functional equation, since the gamma arguments stay
positive throughout the required negative range.
The stronger
\texttt{ThmBop}/\texttt{ThmAop} and
\texttt{PartsE}/\texttt{PartsW} include positive scalar operator
identities and spanning assertions; none is inferred here from
rationality or finite support. The source ledger below locates
actual producers for these tasks, and records which smaller
arithmetic subclaims have been discharged in this section.
